\chapter{Experimento e Coleta de Dados}\label{chp:exp}

Esta seção descreve o experimento e o processo de coleta de dados..... 

\section{Descrição do Experimento}

\section{Processo de coleta de dados}

\section{Resultados Obtidos}

\textbf{ATENÇÃO}
%
%SSDDSDFSFDF
%SFSDFDFFSF
%SFFFSDFF
%
%

\section{Survey}

Nesta etapa, foi aplicado um \textit{survey} que foi elaborado questões semi-estruturadas. Algumas destas questões, ficaram em aberto, de forma proposital para permitir que os entrevistados pudessem falar sobre o assunto da pesquisa, de uma forma mais livre, sem ficar limitados às alternativas das questões. 

O \textit{survey} seguiu as recomendações de análise de Kitchenham e Pfleeger~\cite{pfleeger2001principles}.

O objetivo de pesquisa do estudo foi:

\begin{enumerate}
    \item q1
    \item q2
\end{enumerate}


















...autor, ano de publicação, título do trabalho, objetivo, metodologia aplicada, principais resultados (Tabela 1).

\begin{table}[htbp] 
\centering
\caption{Caracterização das publicações selecionadas} \label{tabela-frameworks}
\begin{tabular}{|p{4cm}|p{4cm}|p{4cm}|p{4cm}|}
\hline
\textbf{TÍTULO DO TRABALHO AUTOR ANO DE PUBLICAÇÃO} & \textbf{OBJETIVO} & \textbf{METODOLOGIA} & \textbf{PRINCIPAIS RESULTADOS}\\ \hline 

Titulo
&
Objetivo
&
Estudo de campo.
&
Resultado.
\\ \hline	
Titulo

&
Objetivo

& 
Revisão de
literatura.

& 
Resultado.

\\ \hline
Titulo

& 
Objetivo

& 
Estudo de campo.
& 
Resultado.

\\ \hline
Titulo

& 
Objetivo

& 
Estudo de campo. 
& 
Resultado.

\\ \hline
Titulo
& 
Objetivo
& 
Estudo de campo. 
& 
Resultado.
\\ \hline
Titulo

& 
Objetivo
& 
Estudo de campo. 
& 
Resultado.
\\ \hline
\end{tabular}
\end{table}

Fonte:dados da pesquisa (2022).

